\documentclass[a4paper,11pt]{article}

\usepackage[utf8]{inputenc}
\usepackage[T1]{fontenc}
\usepackage[german]{babel}
\usepackage[a4paper,margin=2.5cm]{geometry}
\usepackage{enumitem}
\usepackage{titlesec}
\usepackage{longtable}
\usepackage{setspace}
\usepackage{parskip}

\begin{document}

% ---------------------------
% Deckblatt
% ---------------------------

\begin{titlepage}
    \centering

    \vspace*{3cm}

    {\Huge \textbf{Meetingprotokoll}}\\[0.5cm]
    {\Large Statusbericht zur Projektarbeit}\\[2cm]

    \rule{12cm}{0.5pt}\\[0.5cm]

    {\Large Projektarbeit}\\[0.3cm]
    {\large 06. August 2026}

    \rule{12cm}{0.5pt}

    \vfill

    {\large THWS / MINOVA}

\end{titlepage}

% ---------------------------
% Inhaltsverzeichnis
% ---------------------------

\tableofcontents
\newpage

% ---------------------------
% Inhalt
% ---------------------------

\section{Allgemeine Informationen}

\begin{tabular}{ll}
\textbf{Datum:} & 06.08.2026 \\
\textbf{Art des Meetings:} & Projektmeeting / Statusbesprechung \\
\textbf{Nächstes Projektmeeting:} & 21.08.2026 um 16:00 Uhr (Deutschland), 17:00 Uhr (Finnland) \\
\textbf{Abgabe Projektdokumentation:} & 16.09.2026 \\
\textbf{Abschlusspräsentation:} & 23.09.2026 um 17:00 Uhr \\
\end{tabular}

\vspace{1cm}

\section{Projektstatus}

Im Rahmen des Projektmeetings wurde der aktuelle Stand des Avisierungsdashboards vorgestellt. Dabei wurden insbesondere die Filterfunktionen sowie die Darstellung der Informationen und Kommentare im Dashboard betrachtet. Zusätzlich wurden offene fachliche Fragen zur weiteren Gestaltung des Dashboards und zur Integration der verschiedenen Kommunikationskanäle besprochen.

\vspace{0.5cm}

\section{Besprochene Punkte}

Während des Meetings wurden folgende Themen behandelt:

\begin{itemize}[leftmargin=1.2cm]

\item Das aktuelle Avisierungsdashboard wurde vorgestellt und gemeinsam betrachtet. Dabei wurden insbesondere die bereits implementierten Filterfunktionen überprüft.

\item Frau Professor Saueressig merkte an, dass die Kommentare direkt auf der Dashboard-Seite sichtbar sein sollten. Dies war bereits im Rahmen der weiteren Entwicklung des Dashboards vorgesehen.

\item Das Projektteam soll prüfen, ob die Anzeige des Namens des Disponenten für den vorgesehenen Anwendungsfall erforderlich ist.

\item Für die verschiedenen Kommunikationskanäle soll eine einheitliche Nachricht verwendet werden. Dadurch soll sichergestellt werden, dass die übermittelten Informationen unabhängig vom verwendeten Kommunikationskanal konsistent sind.

\end{itemize}

\vspace{0.5cm}

\section{Nächste Schritte}

\begin{itemize}[leftmargin=1.2cm]

\item Im Projektteam prüfen, ob die Anzeige des Namens des Disponenten für den vorgesehenen Anwendungsfall erforderlich ist.

\item Eine einheitliche Nachricht für die verschiedenen Kommunikationskanäle definieren und bei der weiteren Integration berücksichtigen.

\end{itemize}

\vspace{0.5cm}

\section{Nächste Projektphase}

Die nächste Projektphase konzentriert sich auf die Weiterentwicklung des Avisierungsdashboards sowie auf die Integration der verschiedenen Kommunikationskanäle. Dabei soll insbesondere eine einheitliche Nachrichtenstruktur für alle Kommunikationskanäle umgesetzt werden.

\vspace{0.5cm}

\section{Nächste Meilensteine}

\begin{itemize}[leftmargin=1.2cm]

\item Weiterentwicklung des Avisierungsdashboards.

\item Umsetzung einer einheitlichen Nachricht für die verschiedenen Kommunikationskanäle.

\end{itemize}

\vspace{1cm}

\section{Zusammenfassung}

Im Projektmeeting wurde der aktuelle Stand des Avisierungsdashboards vorgestellt und insbesondere die Filterfunktion betrachtet. Frau Professor Saueressig merkte an, dass die Kommentare direkt auf der Dashboard-Seite sichtbar sein sollten. Dies war bereits im Rahmen der weiteren Entwicklung des Dashboards vorgesehen. Darüber hinaus soll das Projektteam prüfen, ob die Anzeige des Namens des Disponenten für den vorgesehenen Anwendungsfall erforderlich ist. Für die verschiedenen Kommunikationskanäle soll eine einheitliche Nachricht verwendet werden, um eine konsistente Kommunikation sicherzustellen. Das nächste Projektmeeting findet am 21.08.2026 um 16:00 Uhr deutscher Zeit bzw. 17:00 Uhr finnischer Zeit statt.
\end{document}